\documentclass[12pt,a4paper]{article}
\usepackage{graphicx}
\usepackage{amsmath}
\usepackage{amssymb}
\usepackage{amsthm}
\usepackage{url}
\usepackage{booktabs}
\usepackage{hyperref}
\usepackage{setspace}
\usepackage{geometry}
\usepackage{background}
\usepackage{titlesec}
\usepackage{tabularx}
\usetikzlibrary{calc}
\usepackage{subcaption}
\usepackage{datetime}
\usepackage{lscape}
\usepackage{longtable}
\usepackage{float}
\usepackage{algorithm}
\usepackage{algorithmic}
\usepackage{listings}
\usepackage{xcolor}
\usepackage{adjustbox}
\usepackage{rotating}

\definecolor{codegray}{rgb}{0.5,0.5,0.5}
\definecolor{codegreen}{rgb}{0,0.6,0}
\definecolor{codeblue}{rgb}{0,0,0.8}
\definecolor{backcolor}{rgb}{0.95,0.95,0.92}

\hypersetup{
    colorlinks=true,
    citecolor=blue,
    linkcolor=blue,
    urlcolor=blue
}

\lstdefinestyle{mystyle}{
    backgroundcolor=\color{backcolor},   
    commentstyle=\color{codegreen},
    keywordstyle=\color{codeblue},
    numberstyle=\tiny\color{codegray},
    stringstyle=\color{codegreen},
    basicstyle=\ttfamily\footnotesize,
    breakatwhitespace=false,         
    breaklines=true,                 
    captionpos=b,                    
    keepspaces=true,                 
    numbers=left,                    
    numbersep=5pt,                  
    showspaces=false,                
    showstringspaces=false,
    showtabs=false,                  
    tabsize=2
}

\lstset{style=mystyle}

\renewcommand{\algorithmicrequire}{\textbf{Input:}}
\renewcommand{\algorithmicensure}{\textbf{Output:}}
\renewcommand{\algorithmiccomment}[1]{\hfill{\small\textit{#1}}}

\geometry{
 a4paper,
 total={170mm,257mm},
 left=20mm,
 top=20mm,
}

\backgroundsetup{angle=0, scale=1, vshift=-5ex,
  contents={\tikz[overlay, remember picture]
    \draw [rounded corners=7pt, line width=1.8pt,
           color=black, fill=white!50, double=black!120]
           ($(current page.north west)+(0.4cm,-0.4cm)$)
           rectangle ($(current page.south east)+(-0.4,0.4)$);}}

\titleformat{\chapter}[display]
  {\normalfont\bfseries}
  {\makebox[\textwidth][l]{\fontsize{36}{42}\selectfont\chaptername\ \thechapter}}
  {0pt}
  {\centering\fontsize{28}{34}\selectfont}
  []

\titlespacing*{\chapter}{0pt}{0pt}{14pt}

\titleformat{\section}
  {\normalfont\fontsize{12}{14}\bfseries}{\thesection}{1em}{}
\titlespacing*{\section}{0pt}{*2}{*1.5}

\titleformat{\subsection}
  {\normalfont\fontsize{12}{14}\bfseries\itshape}{\thesubsection}{1em}{}
\titlespacing*{\subsection}{0pt}{*1.5}{*1}

\titleformat{\subsubsection}
  {\normalfont\fontsize{12}{14}\bfseries\itshape}{\thesubsubsection}{1em}{}
\titlespacing*{\subsubsection}{0pt}{*1}{*0.5}

\usepackage{fancyhdr}
\fancypagestyle{plain}{
  \fancyhf{}
  \renewcommand{\headrulewidth}{0pt}
  \renewcommand{\footrulewidth}{0pt}
}

\fancypagestyle{chapterstyle}{
  \fancyhf{}
  \fancyhead[L]{\leftmark}
  \fancyhead[R]{\thepage}
  \fancyfoot[C]{\small B.M.S College of Engineering, Bangalore}
  \renewcommand{\headrulewidth}{1.5pt}
  \renewcommand{\footrulewidth}{1.0pt}
}

\begin{document}
\linespread{1.39}\selectfont
\pagestyle{plain}
\onehalfspacing

\begin{titlepage}
    \begin{center}
        \textsc{\Large \textbf{\textcolor{blue}{B.M.S. COLLEGE OF ENGINEERING}}}\\[0.2cm]
        \textsc{\large \textbf{(Autonomous College under VTU)}}\\[0.2cm]
        \textsc{\large \textbf{Bull Temple Road, Basavangudi, Bangalore - 560019}}\\[0.5cm]
        
        \includegraphics[scale=0.75]{logo.png}\\[0.2cm]
        \textit{A report on}\\[0.2cm]
        \textsc{\textbf{\textcolor{red}{Cloud Computing - AAT}}}\\[0.2cm]
        \textit{in fulfillment of the requirements for the award of degree}\\[0.2cm]
        
        \textbf{BACHELOR OF ENGINEERING}\\[0.15cm]
        \textit{in}\\[0.2cm]
        \textbf{Computer Science \& Engineering (Data Science)}\\[0.2cm]
        
        \textit{ by}\\[0.2cm]
         \textbf{{Athish Raj (1BM23CD009)}}\\[0.2cm]

        \textbf{Under the Guidance of}\\[0.25cm]
        \textbf{Dr. Professor Name}\\[0.15cm]
        Assistant Professor\\[0.55cm]
        \textbf{Department of CSE (Data Science)}, \\
        \textbf{B.M.S. College of Engineering}\\[0.3cm]

        \textbf{AY 2025-2026}
    \end{center}
\end{titlepage}

%---------------------------------------------------------------------
\section{Abstract}
This project demonstrates the design, deployment, and cloud-native scaling of a high-performance Data Science personal portfolio website using Amazon Web Services (AWS). Built with a modern HTML/CSS standalone architecture featuring a dark "glassmorphism" aesthetic, the application is hosted dynamically to exhibit core cloud computing concepts. By leveraging AWS Application Load Balancers (ALB) and Auto Scaling Groups (ASG) across multiple EC2 instances running Nginx on Amazon Linux 2023, the architecture natively ensures high availability, horizontal scalability, and robust fault tolerance.

\medskip
\noindent\textbf{Live Website URL:} \url{http://13.127.216.228}

%---------------------------------------------------------------------
\section{Introduction}
The primary objective of this Alternative Assessment Tool (AAT) project is to practically apply foundational cloud computing principles by deploying a production-ready web application. 
\begin{itemize}
    \item \textbf{Purpose:} To transition a local frontend portfolio into a highly available cloud architecture, migrating away from inherently limited static hosting (e.g., GitHub Pages) toward an enterprise-grade infrastructure.
    \item \textbf{Concepts Demonstrated:} Load Balancing (distributing web traffic evenly), Horizontal Scaling (dynamic instance allocation), and Fault Tolerance (auto-healing crashed servers).
    \item \textbf{Provider Justification:} AWS was selected due to its status as the industry standard for Infrastructure as a Service (IaaS). The deployment extensively utilized Free Tier eligible configurations (t2.micro/t3.micro instances) while successfully satisfying all enterprise networking constraints.
\end{itemize}

%---------------------------------------------------------------------
\section{Architecture Overview}
The cloud environment was designed around a multi-tier, fault-tolerant infrastructure capable of handling localized data center outages seamlessly without impacting end-user experience.

\begin{figure}[H]
    \centering
    % Replace with your diagram
    \includegraphics[width=0.55\linewidth]{figures/architecture.png}
    \caption{High-level deployment architecture representing User Traffic parsed by the Load Balancer to active ASG targets.}
    \label{fig:arch}
\end{figure}

The architecture relies on the following key components:
\begin{itemize}
    \item \textbf{Application Load Balancer (ALB):} Acts as the primary ingress point. It actively pings target EC2 instances to verify health status and distributes incoming HTTP traffic across geographically distinct Availability Zones.
    \item \textbf{Auto Scaling Group (ASG):} Constantly monitors the target group. If traffic spikes or an instance crashes, the ASG automatically terminates the dead node and provisions a fresh VM clone from the established Launch Template.
    \item \textbf{EC2 Instances (VMs):} The core computational nodes running Amazon Linux 2023. Governed by a custom Launch Template, they host the compiled HTML Application securely using the Nginx Web Server.
    \item \textbf{Dynamic Metadata Endpoint:} A localized script pulls the node's unique AWS Instance ID directly from the hypervisor layer (sysfs hardware files), creating a pseudo-API text file that the frontend consumes to prove load distribution in real-time.
\end{itemize}

%---------------------------------------------------------------------
\section{Website Design and Implementation}
The portfolio is bundled into a single standalone \texttt{index.html} payload, minimizing dependency trees and eliminating backend operational overhead. It utilizes an advanced \textit{Glassmorphism} UI built with custom CSS variables, offering responsive dynamic layout transitions and built-in JavaScript Dark/Light theme toggles. To visualize load balancing, an asynchronous \texttt{fetch()} logic is embedded at runtime to request server-side identification endpoints.

\begin{figure}[H]
    \centering
    \begin{subfigure}{0.45\linewidth}
        \includegraphics[width=0.45\linewidth]{figures/site_home.png}
        \caption{Home page}
    \end{subfigure}
    \hfill
    \begin{subfigure}{0.45\linewidth}
        \includegraphics[width=0.45\linewidth]{figures/site_projects.png}
        \caption{Projects page}
    \end{subfigure}
    \caption{Screenshots displaying the responsive Glassmorphism design and the dynamic "Served by Instance" banner.}
\end{figure}

%---------------------------------------------------------------------
\section{Cloud Setup Process}

\subsection{Provider and Service Selection}
The system leverages AWS EC2 for raw compute, AWS Target Groups for networking ingress, and AWS IMDSv2 metadata logic for API spoofing during deployment.

\subsection{Custom AMI and Instance Configuration}
To bypass standard remote code repository restrictions (e.g. GitHub/S3 dependencies), a Custom Amazon Machine Image (AMI) strategy was adopted. A single "Golden Instance" was provisioned with all application code baked directly into its local drive. This perfectly configured server image served as the strict blueprint for all newly allocated VMs spun up by the Auto Scaling logic.

\begin{figure}[H]
    \centering
    \includegraphics[width=0.45\linewidth]{figures/instance_creation.png}
    \caption{Creating the compute instance on the cloud console.}
\end{figure}
\subsection{Deployment Steps (Nginx Setup)}
Commands used to bootstrap the application to the high-performance internet-facing Linux daemon:

\begin{lstlisting}[language=bash,caption={Deploying the bundled architecture on Amazon Linux 2023}]
# Create a secure copy transaction using the AWS SSH Key
scp -i portfolio-key.pem index.html ec2-user@<public-ip>:/home/ec2-user/

# Move dependencies to the public-facing Nginx web directory
sudo cp /home/ec2-user/index.html /usr/share/nginx/html/

# Extract the AWS Instance hardware ID and serve it as an API Endpoint
sudo sh -c 'cat /sys/devices/virtual/dmi/id/board_asset_tag > /usr/share/nginx/html/instance'

# Reload systemd web service without dropping current connections
sudo systemctl reload nginx
\end{lstlisting}

\begin{figure}[H]
    \centering
    \includegraphics[width=0.45\linewidth]{figures/deployment.png}
    \caption{Website running successfully on the cloud instance.}
\end{figure}

%---------------------------------------------------------------------
\section{Demonstration of Cloud Concepts}

\subsection{Load Balancing}
The AWS Application Load Balancer is bound to an HTTP Target Group spanning at least two distinct AZs (e.g., \texttt{ap-south-1a} and \texttt{ap-south-1b}). The ALB resolves ingress using an intelligent Round Robin paradigm. 
\begin{itemize}
    \item \textbf{Visual Verification:} The Javascript dynamically pulls the endpoint. On hard page refreshes, the banner cleanly alternates between Instance IDs (e.g., \texttt{i-0abcdef} vs \texttt{i-0123456}), physically proving that the Load Balancer has successfully routed the user to a different cluster server.
\end{itemize}

\begin{figure}[H]
    \centering
    \includegraphics[width=0.45\linewidth]{figures/load_balancer.png}
    \caption{Load balancer configuration and healthy targets.}
\end{figure}

\subsection{Fault Tolerance \& Auto-Healing}
Resilience is handled autonomously by the configured ASG Health check protocols. 
\begin{itemize}
    \item If a single underlying EC2 backend is stopped manually via the AWS Console, the ALB immediately detects a timeout.
    \item Traffic is mathematically routed specifically away from the failing node to the surviving instance. 
    \item Within seconds, the ASG invokes its Minimum Capacity directive (set to 2). It terminates the irrecoverable node and boots a new, fully operative clone utilizing the baked AMI blueprint.
\end{itemize}

\begin{figure}[H]
    \centering
    \includegraphics[width=0.45\linewidth]{figures/fault_tolerance.png}
    \caption{Website remains available after one instance is stopped.}
\end{figure}

%---------------------------------------------------------------------
\section{Testing and Results}
Deployment was verified consistently across the application layers.

\begin{table}[H]
\centering
\begin{tabular}{@{}lll@{}}
\toprule
\textbf{Test} & \textbf{Expected Result} & \textbf{Observed Result} \\
\midrule
Nginx HTTP Ingress    & 200 OK Response   & 200 OK \\
ALB Traffic Splitting & Alternating IDs   & Banner swapped instances on refresh \\
Forced VM Failure     & Zero User Downtime& Seamlessly failed over to secondary node \\
\bottomrule
\end{tabular}
\caption{System Resilience Testing Summary}
\end{table}

%---------------------------------------------------------------------
\section{Challenges and Learnings}
A significant operational hurdle encountered was related to the stringent AWS IMDSv2 security policies deployed by default on modern Amazon Linux 2023 kernels. Standard \texttt{curl} HTTP requests to fetch instance-metadata dynamically over port 80 repeatedly returned `401 Unauthorized` errors. 

To resolve this without writing convoluted secondary token-authentication scripts, extensive system research led to querying the hypervisor's directly mapped SysFS directory (\texttt{/sys/devices/virtual/dmi/id/board\_asset\_tag}). This provided a stable, fool-proof physical read operation that allowed the Nginx server to securely mock backend API capabilities seamlessly.

%---------------------------------------------------------------------
\section{Conclusion}
This deployment successfully transitioned a statically designed frontend interface into a fully scalable cloud ecosystem. Core IaaS concepts including Application Load Balancing, high-availability zone partitioning, and self-healing auto-scaling arrays were thoroughly modeled and verified. The workflow effectively simulated an enterprise-ready pipeline while establishing stable groundwork for more complicated structural implementations in the future, such as Kubernetes containerization or CI/CD integrations.

%---------------------------------------------------------------------
\end{document}
